\documentclass[12pt,a4paper]{article}

\usepackage[T1]{fontenc}
\usepackage[utf8]{inputenc}
\usepackage[ngerman]{babel}
\usepackage{lmodern}
\usepackage{microtype}

\usepackage{amsmath,amssymb,amsthm,mathtools}
\usepackage[a4paper,margin=2.4cm]{geometry}
\usepackage{booktabs}
\usepackage{xcolor}
\usepackage[hidelinks,unicode]{hyperref}
\usepackage[most]{tcolorbox}
\usepackage{enumitem}
\usepackage{listings}

\definecolor{bewiesengruen}{RGB}{30,100,50}
\definecolor{warnorange}{RGB}{200,100,10}
\definecolor{lueckenrot}{RGB}{180,40,40}
\definecolor{kernblau}{RGB}{20,60,130}

\tcbset{
  bewiesen/.style={colback=green!7,colframe=bewiesengruen,fonttitle=\bfseries},
  heuristik/.style={colback=orange!10,colframe=warnorange,fonttitle=\bfseries},
  offen/.style={colback=red!8,colframe=lueckenrot,fonttitle=\bfseries},
  kern/.style={colback=blue!5,colframe=kernblau,fonttitle=\bfseries},
}

\newtheorem{satz}{Satz}[section]
\newtheorem{lemma}[satz]{Lemma}
\newtheorem{definition}[satz]{Definition}
\newtheorem{bemerkung}[satz]{Bemerkung}
\newtheorem{vermutung}[satz]{Vermutung}

\newcommand{\vii}{\nu_2}
\newcommand{\N}{\mathbb{N}}
\newcommand{\Z}{\mathbb{Z}}
\newcommand{\Ztwo}{\mathbb{Z}_2}
\newcommand{\Ztwost}{\mathbb{Z}_2^{\ast}}
\newcommand{\U}{\mathcal{U}}
\newcommand{\Col}{\mathrm{Col}}
\newcommand{\disttwo}{\operatorname{dist}_2}

\lstset{
  basicstyle=\ttfamily\scriptsize,
  breaklines=true,
  frame=single,
  numbers=left,
  numberstyle=\tiny,
  language=,
  showstringspaces=false
}

\title{%
  Die Ausnahmemenge $E\subset\mathbb{Z}_2$ und die 2-adische Distanz\\
  $\mathrm{dist}_2(n,E)$ als Uniformitäts-Objekt
}
\author{Thomas Hoffbauer}
\date{16.\ Juni 2026}

\begin{document}
\maketitle

\begin{abstract}
Die Collatz-Vermutung wird hier nicht erneut in der klassischen Form
``jede Bahn erreicht $1$?'' angegangen.
Stattdessen fragen wir, warum keine Trajektorie dauerhaft vom
2-adischen Attraktorbereich der triviale Zyklus-Nachbarschaft getrennt bleibt.
Wir definieren die odd-to-odd-Abbildung $U$ auf $\Ztwost$,
präzisieren eine Ausnahmemenge $E\subset\Ztwo$ als 2-adischen Limes
nicht-trivialer Bahnen und studieren $\disttwo(n,E)$.
Bekannte Widerlegungen falscher Attraktoren (Abstand zu $1$, naive
$c\cdot 2^{-\log n}$-Schranke) werden referenziert.
Die verbleibende Uniformitäts-Vermutung wird punktweise formuliert.
Ein Lean-Gerüst \texttt{collatz\_z2\_attraktor.lean} skizziert die Formalisierung.
\end{abstract}

\tableofcontents
\newpage

% =================================================
\section{Motivation: vom globalen Ziel zum 2-adischen Attraktor}
\label{sec:motivation}
% =================================================

Die klassische Collatz-Frage lautet: Konvergiert jede Bahn $n,\Col(n),\Col^2(n),\ldots$
nach endlich vielen Schritten zum Zyklus $(1,2,4)$?
Bewiesen ist eine reiche modulare Struktur (EABC, C-Ketten-Endlichkeit,
LTE-Dichte, negativer mittlerer Drift in Blockmodellen), aber nicht die
punktweise Aussage f\"ur alle $n\in\N$.

\begin{tcolorbox}[kern,title={Verschiebung der Leitfrage}]
\emph{Benutzer-Leitfrage:} Nicht prim\"ar ``erreicht jede Bahn $1$?'',
sondern: \textbf{Warum kann keine Trajektorie dauerhaft vom 2-adischen
Attraktorbereich der trivialen Nachbarschaft getrennt bleiben?}
\end{tcolorbox}

Intuitiv tr\"agt die odd-to-odd-Dynamik
\[
  U(n)=\frac{3n+1}{2^{\vii(3n+1)}},\qquad n\ \text{ungerade},
\]
Information in $\Ztwo$ nach links (h\"ohere $2$-Potenzen in Differenzen),
w\"ahrend der bekannte Attraktor $1,2,4,\ldots$ eine \emph{2-adische Umgebung}
von $1$ beschreibt.
Die offene Uniformit\"atsl\"ucke ist dann: Aus fast-sicherem oder
mittlerem Verhalten folgt nicht automatisch, dass ein \emph{festes} $n$
nicht in einer pathologischen 2-adischen Komponente verharrt.
Dieses Dokument fixiert die Begriffe $E$ und $\disttwo(n,E)$ als
gemeinsame Sprache f\"ur diesen Angriff.

% =================================================
\section{Definitionen}
\label{sec:def}
% =================================================

\subsection{Odd-to-odd Collatz auf $\Ztwost$}

Sei $\Ztwo$ der Ring der 2-adischen ganzen Zahlen und
$\Ztwost=\Ztwo\setminus\{0\}$.
Jedes ungerade $n\in\N$ identifizieren wir mit seinem Bild in $\Ztwo$.

\begin{definition}[Odd-to-odd-Abbildung]
\label{def:U}
F\"ur ungerades $x\in\Ztwo$ mit $3x+1\neq 0$ setzen wir
\[
  U(x):=\frac{3x+1}{2^{\vii(3x+1)}}\in\Ztwo,
\]
sofern die $2$-Bewertung $\vii(3x+1)$ endlich ist (f\"ur alle $x\in\Ztwost$ der Fall).
Einschr\"ankung auf $\N$ liefert die \"ubliche Collatz-Kette auf ungeraden Startwerten.
\end{definition}

\begin{bemerkung}
Die Abbildung $U$ ist auf $\Ztwost$ wohldefiniert und kontinuierlich in der
2-adischen Topologie; f\"ur die punktweise Collatz-Vermutung reicht jedoch
die Einschr\"ankung $U^K:\N\to\N$ nicht aus --- genau hier liegt die L\"ucke.
\end{bemerkung}

\subsection{Trivialer Attraktor und $E$}

\begin{definition}[Triviale 2-adische Nachbarschaft]
\[
  A_{\mathrm{triv}}:=\{x\in\Ztwo : |x-1|_2 < 1\}
  = 1+2\Ztwo .
\]
Entspricht Bahnen, die in der Collatz-Dynamik in den Zyklus $(1,2,4)$ fallen
(nach Einbettung ungerader Schritte in die volle Dynamik).
\end{definition}

\begin{definition}[Ausnahmemenge $E$ (Limes-Version)]
\label{def:E}
Eine Trajektorie $(x_k)_{k\geq 0}$ mit $x_{k+1}=U(x_k)$ hei\ss e
\emph{trivial}, wenn es $K$ gibt mit $x_K\in A_{\mathrm{triv}}$.
Sonst hei\ss e sie \emph{ausnahme-typisch} (divergent im 2-adischen Sinne
oder nicht-trivial zyklisch).
F\"ur $K\in\N$ sei $E_K\subset\N$ die Menge der $n\leq K$, deren
\emph{nat\"urliche} Collatz-Bahn ausnahme-typisch ist (divergiert oder
Zyklus $\neq(1,2,4)$).
Wir setzen
\[
  E:=\bigcap_{K\in\N}\overline{\{x\in\Ztwo : \exists n\in E_K,\ |x-n|_2\leq 2^{-K}\}}^{\;2\text{-adisch}},
\]
also den Schnitt der abgeschlossenen H\"ullen divergenter bzw.\ nicht-trivial
zyklischer Pr\"afixe in $\Ztwo$.
\end{definition}

\begin{tcolorbox}[heuristik,title={Alternative Lesarten von $E$}]
\begin{enumerate}[label=(\alph*),nosep]
  \item \textbf{Divergenz in $\N$:} $n$ mit unbeschr\"ankter Bahn --- klassisch,
    aber ohne Beweis der Leerheit.
  \item \textbf{Nicht-triviale Zyklen:} endliche Perioden $\neq(1,2,4)$;
    bekannt unentschieden, in $\Ztwo$ als endliche Orbits modul $2^M$ testbar.
  \item \textbf{Haar-Null vs.\ $\N$:} In $\Ztwo$ ist $A_{\mathrm{triv}}$ eine
    offene Umgebung; ``fast alle'' 2-adische Punkte sind trivial \emph{als Ma\ss e},
    nicht als nat\"urliche Zahlen.
  \item \textbf{Pr\"afix-Limes:} $E$ als oben --- betont 2-adische Koh\"arenz
    langsam wachsender Pr\"afixe statt einzelner Diophantischer Ausnahmen.
\end{enumerate}
Dieses Dokument verwendet Definition~\ref{def:E}; die Alternativen sind
\"aquivalent \emph{nur} unter zus\"atzlichen, derzeit unbewiesenen Annahmen.
\end{tcolorbox}


\subsection{Beispiel: LTE-Familie und Abstand zu $1$}

Für $n_0=2^{k+1}3^r-1$ mit $k,r\geq 1$ gilt $\vii(n_0+1)=k+1$, also
$|n_0-(-1)|_2=2^{-(k+1)}$ und insbesondere für $k=1$ oft $\disttwo(n_0,1)=1/2$.
Diese Familie illustriert, warum Uniformitätsaussagen über $\disttwo(\cdot,1)$
nicht den Attraktor $A_{\mathrm{triv}}$ treffen (vgl.\ PR~\#17).

\begin{center}
\begin{tabular}{@{}lll@{}}
\toprule
Start $n_0$ & $\vii(n_0-1)$ & $\disttwo(n_0,1)$ \\
\midrule
$4\cdot 3^r-1$ & $\geq 2$ & $\leq 1/4$ oder $1/2$ (je nach $r$) \\
$2^{k+1}-1$ (Mersenne-artig) & $k+1$ & $2^{-k}$ \\
\bottomrule
\end{tabular}
\end{center}

\subsection{Verbindung zu $E_K$ und natürlichen Zahlen}

Für festes $K$ ist $E_K$ endlich und berechenbar, sofern man Collatz-Bahnen
bis zur Divergenzsschwelle oder Zykluserkennung simuliert.
Die Uniformitäts-Vermutung verlangt, dass kein $n\in\N$ für alle $K$
in $E_K$ verbleibt --- eine punktweise Aussage, die strenger ist als
``$E_K$ hat Dichte $0$ für jedes $K$''.


\subsection{2-adische Metrik und $\disttwo(n,E)$}

F\"ur $x,y\in\Ztwo$, $x\neq y$, setzen wir
\[
  |x-y|_2 := 2^{-\vii(x-y)} .
\]
F\"ur $n\in\N$ und $E\subset\Ztwo$
\[
  \disttwo(n,E):=\inf_{e\in E}|n-e|_2 ,
\]
mit der Konvention $\disttwo(n,\emptyset)=1$, falls $E=\emptyset$.

\begin{definition}[Endliche Approximation $E_K$]
F\"ur Rechnungen in $\N$ nutzen wir parallel
\[
  E_K:=\{n\in\N : n\leq K\ \text{und Bahn ausnahme-typisch}\},
\]
und die uniforme Bedingung $U^K(n)\notin E_K$ aus \S\ref{sec:uniform}.
\end{definition}

\begin{bemerkung}
F\"ur $a,b\in\Z$ gilt $\disttwo(a,b)=2^{-\vii(a-b)}$.
Die Distanz zu $1$ ist \emph{nicht} dasselbe wie $\disttwo(n,E)$; siehe \S\ref{sec:widerlegt}.
\end{bemerkung}

% =================================================
\section{Was widerlegt wurde}
\label{sec:widerlegt}
% =================================================

Aus \texttt{collatz\_angriff\_uniformitaet.tex} (Pull Request \#17) und numerischen
Gegenbeispielen folgt:

\begin{tcolorbox}[bewiesen,title={Falscher Attraktor: Abstand zu $1$}]
Die naive Uniformit\"atsform
\[
  \disttwo(n,1)\geq c\cdot 2^{-\lfloor\log_2 n\rfloor}
\]
ist \textbf{widerlegt}: F\"ur LTE-Familien $n_0=2^{k+1}3^r-1$ bleibt
$\disttwo(n_0,1)=2^{-\vii(n_0-1)}$ oft bei $1/2$ konstant, w\"ahrend
$c\cdot 2^{-\lfloor\log_2 n\rfloor}\to 0$.
Der relevante Attraktor ist $A_{\mathrm{triv}}$ bzw.\ die Menge $E$,
nicht der einzelne Punkt $1\in\Ztwo$.
\end{tcolorbox}

\begin{tcolorbox}[bewiesen,title={Widerlegte $c\cdot 2^{-\log n}$-Form}]
Ebenfalls widerlegt (Proposition in PR~\#17): Jede Schranke der Form
$\disttwo(n,E)\geq c\cdot 2^{-\lfloor\log_2 n\rfloor}$ mit festem $c>0$,
die nur $\disttwo(\cdot,1)$ testet, scheitert an expliziten $n_0$.
Eine korrekte Uniformit\"atsaussage muss $E$ (oder $E_K$) explizit tragen
und darf den Attraktor nicht durch $1$ ersetzen.
\end{tcolorbox}

Diese Negativresultate sind wertvoll: Sie verbieten eine Klasse scheinbar
plausibler 2-adischer Heuristiken und zwingen zur pr\"azisen Definition~\ref{def:E}.

% =================================================
\section{Uniformit\"ats-Vermutung (pr\"azisiert)}
\label{sec:uniform}
% =================================================

\begin{vermutung}[Punktweise Uniformit\"at / Austritt aus $E_K$]
\label{verm:uniform}
F\"ur jedes $n\in\N$ existiert $K(n)<\infty$, sodass die odd-to-odd-Trajektorie
nach h\"ochstens $K(n)$ Schritten \emph{nicht} in der endlichen Ausnahme-Approximation
$E_{K(n)}$ liegt; \"aquivalent (unter Collatz-Konvergenz): $U^{K(n)}(n)\notin E_{K(n)}$.
\end{vermutung}

\begin{tcolorbox}[offen,title={\"Aquivalente Distanz-Formulierung (bedingt)}]
\emph{Wenn} $E\subset\Ztwo$ die in Definition~\ref{def:E} beschriebene
Ausnahmemenge ist und \emph{wenn} jede nicht-triviale Bahn positiven Abstand
zu $A_{\mathrm{triv}}$ hat, dann ist Vermutung~\ref{verm:uniform} \"aquivalent zu:
\[
  \forall n\in\N\ \exists K(n):\quad \disttwo(U^{K(n)}(n),E)\geq 2^{-\alpha}
\]
f\"ur eine feste Tiefenschranke $\alpha$ (noch zu kalibrieren), oder schw\"acher:
$\disttwo(n,E)$ kann nicht entlang der ganzen Bahn gegen $0$ tendieren, ohne
$A_{\mathrm{triv}}$ zu treffen.
\end{tcolorbox}

\begin{tcolorbox}[kern]
Die Collatz-Vermutung impliziert $E\cap\N=\emptyset$ und damit
$\disttwo(n,E)=1$ f\"ur alle $n$; umgekehrt folgt aus Vermutung~\ref{verm:uniform}
allein ohne Zusatzinformation \emph{nicht} automatisch Konvergenz nach $1$.
Die Uniformit\"ats-Vermutung ist eine \textbf{Zwischenstufe}: strukturell
schw\"acher als voller Collatz, aber st\"arker als fast-\"uberall-Aussagen
(Tao, Birkhoff auf Blockmodellen).
\end{tcolorbox}

% =================================================
\section{Lean-Skizze: \texttt{collatz\_z2\_attraktor.lean}}
\label{sec:lean}
% =================================================

Das Lean-Ger\"ust importiert Padics aus Mathlib und trennt bewiesene
$\vii$-Lemmas von noch offenen 2-adischen Normen.

\begin{lstlisting}
-- collatz_z2_attraktor.lean (Auszug)
import Mathlib.NumberTheory.Padics.PadicInt
import Mathlib.NumberTheory.Padics.PadicVal

namespace CollatzZ2

def oddPart (n : Nat) : Nat := ...

def U_odd (n : Nat) (h : Odd n) : Nat := ...

-- dist_2 auf Z: |a-b|_2 = 2^{-padicValNat 2 (a-b)}
def dist2Nat (a b : Nat) : Rat := ...

def ExceptionSetApprox (K : Nat) : Finset Nat := ...

theorem dist_to_one_not_uniform :
    Exists (fun n => dist2Nat n 1 = 1/2) := by sorry

end CollatzZ2
\end{lstlisting}

\begin{tcolorbox}[heuristik,title={Mathlib-Roadmap}]
\textbf{Jetzt:} \texttt{padicValNat} f\"ur $\vii(n-m)$ und LTE-Dichte
(vgl.\ \texttt{collatz\_uniformity.lean}).\\
\textbf{N\"achster Schritt:} \texttt{PadicInt.norm} bzw.\ \texttt{Padic.norm}
f\"ur $|x-y|_2$ auf eingebetteten $\N\hookrightarrow\Ztwo$.\\
\textbf{Offen:} Formale Definition von $E$ als Filter-Limes; Mathlib liefert
Topologie auf $\Ztwo$, aber keine Collatz-spezifische Ausnahmemenge.
\texttt{sorry} nur an diesen L\"ucken, nicht bei bereits bewiesenen Valuations-Lemmas.
\end{tcolorbox}

% =================================================
\section{Offene Fragen}
\label{sec:offen}
% =================================================

\begin{enumerate}[label=\textbf{F\arabic*.}]
  \item Ist $E\cap\N=\emptyset$ \"aquivalent zur Collatz-Vermutung, oder strikt schw\"acher?
  \item Gibt es $c>0$ und eine Funktion $\psi(n)\to 0$ (nicht $2^{-\log n}$), sodass
    $\disttwo(n,E)\geq c\cdot\psi(n)$ f\"ur alle bekannten Gegenfamilien gilt?
  \item L\"asst sich Vermutung~\ref{verm:uniform} aus einem endlichen
    mod-$2^M$-Zertifikat plus LTE-Reset (gerades $N$) ableiten?
  \item Welche Menge in $\Ztwo$ ist die kleinste abgeschlossene $U$-invariante
    H\"ulle von $A_{\mathrm{triv}}$ --- und stimmt sie mit dem Limes-Definierten $E^c$ \"uberein?
  \item Kann $\disttwo(n,E)$ f\"ur LTE-Familien $2^{k+1}3^r-1$ asymptotisch
    numerisch stabilisiert werden, um eine vermutete untere Schranke zu kalibrieren?
\end{enumerate}

\begin{tcolorbox}[offen,title={Ehrliche Einordnung}]
Dieses Arbeitspaket definiert Begriffe und sammelt widerlegte Heuristiken;
es beansprucht keinen Beweis der Collatz-Vermutung noch von
Vermutung~\ref{verm:uniform}.
\end{tcolorbox}

\end{document}
